\documentclass[../main.tex]{subfiles}
\begin{document}

\chapter{Checklists}
\label{ch:checklists}

These checklists collect the criteria and rules scattered through the week chapters into a single place for use during the exam. Each box names the section it is drawn from. The square boxes are for ticking against your own work.

\begin{checklistbox}{Research Objective Criteria \quad (\autoref{sec:objectives-rqs})}
\begin{itemize}[leftmargin=2em]
  \item[\(\square\)] \textbf{Informative}: could someone with generic engineering knowledge understand what the project is about?
  \item[\(\square\)] \textbf{Useful}: what is the benefit of the research to the problem or solution?
  \item[\(\square\)] \textbf{Feasible}: is it possible within your capabilities and resources?
  \item[\(\square\)] \textbf{Clear}: is it precise about the project's contribution?
\end{itemize}
\end{checklistbox}

\begin{checklistbox}{Research Question Formulation \quad (\autoref{sec:objectives-rqs})}
\begin{itemize}[leftmargin=2em]
  \item[\(\square\)] The questions are yours to answer, not the client's.
  \item[\(\square\)] There are at least two sub-questions, never only one.
  \item[\(\square\)] Questions are open-ended rather than yes or no.
  \item[\(\square\)] Each question is relevant, researchable, anchored, and precise.
  \item[\(\square\)] The question names the key concept and links to quantitative evidence.
  \item[\(\square\)] The set builds toward a testable main hypothesis.
  \item[\(\square\)] If the technical design struggles, the questions are revisited and edited.
\end{itemize}
\end{checklistbox}

\begin{checklistbox}{Data Quality Checks \quad (\autoref{sec:data-sources})}
\begin{itemize}[leftmargin=2em]
  \item[\(\square\)] Checked for typing errors.
  \item[\(\square\)] Checked for outliers.
  \item[\(\square\)] Checked for errors in coding the data.
  \item[\(\square\)] Checked whether the data meets expectations.
  \item[\(\square\)] A detailed research logbook is kept.
  \item[\(\square\)] A data management plan is in place.
\end{itemize}
\end{checklistbox}

\begin{checklistbox}{Data Management Plan Sections \quad (\autoref{sec:dmp})}
\begin{itemize}[leftmargin=2em]
  \item[\(\square\)] \textbf{I. Data description}: types, formats, collection, purpose, storage, access, and volume.
  \item[\(\square\)] \textbf{II. Documentation and data quality}: methodology, lab notebook, README, and data dictionary.
  \item[\(\square\)] \textbf{III. Storage and backup}: storage location and the 3-2-1 rule of three copies, two media, one off-site.
  \item[\(\square\)] \textbf{IV. Legal and ethical requirements}: confidential data, human subjects, and third-party datasets.
  \item[\(\square\)] \textbf{V. Sharing and long-term preservation}: what is shared, how, and under which license.
  \item[\(\square\)] \textbf{VI. Responsibilities}: lead institution and who owns the data after you leave.
\end{itemize}
\end{checklistbox}

\begin{checklistbox}{Reference Checklist \quad (\autoref{sec:referencing})}
\begin{itemize}[leftmargin=2em]
  \item[\(\square\)] All borrowed illustrations are referenced.
  \item[\(\square\)] All borrowed statements are substantiated with references.
  \item[\(\square\)] Direct citations are in quotation marks and paraphrases are in your own words.
  \item[\(\square\)] Each citation or paraphrase has a correct in-text reference in the chosen style.
  \item[\(\square\)] All references are correct and complete in the chosen style.
  \item[\(\square\)] The list contains all references used in the text.
  \item[\(\square\)] The list contains only references used in the text.
\end{itemize}
\end{checklistbox}

\begin{checklistbox}{Final Writing Checklist \quad (\autoref{sec:writing-style})}
\begin{itemize}[leftmargin=2em]
  \item[\(\square\)] No spelling or grammar errors.
  \item[\(\square\)] Each paragraph conveys one or two ideas.
  \item[\(\square\)] The paragraph format is consistent throughout.
  \item[\(\square\)] The nomenclature is consistent throughout.
  \item[\(\square\)] All acronyms are defined on first use.
  \item[\(\square\)] All figures are readable in PDF.
  \item[\(\square\)] All figures and tables include a caption.
  \item[\(\square\)] Figures are clear without reading the text.
  \item[\(\square\)] The title captures the contents.
  \item[\(\square\)] The introduction captures the ideal reader.
  \item[\(\square\)] Conclusions are understood without reading the rest.
\end{itemize}
\end{checklistbox}

\begin{checklistbox}{Sustainable Research Reflection \quad (\autoref{sec:eop})}
\begin{itemize}[leftmargin=2em]
  \item[\(\square\)] Is there potential value to science or industry from the findings?
  \item[\(\square\)] Can the research be built upon?
  \item[\(\square\)] Is there direct potential harm to the environment, society, or economy as intended?
  \item[\(\square\)] Could the research be used for unintended harmful purposes?
  \item[\(\square\)] Is the desired accuracy realistic given the resources available?
  \item[\(\square\)] Is there a less resource-intensive alternative with reasonable accuracy?
  \item[\(\square\)] Is the sample size justified?
\end{itemize}
\end{checklistbox}

\end{document}
